\documentclass[11pt]{article}

\usepackage[margin=0.9in]{geometry}
\usepackage{amsmath}
\usepackage{booktabs}
\usepackage{graphicx}
\usepackage{float}
\usepackage{hyperref}
\usepackage{array}
\usepackage[table]{xcolor}

\graphicspath{{metric_analysis_files/}}
\sloppy

\definecolor{lightgray}{gray}{0.94}

\newcommand{\simulationassumptions}[4]{%
\begin{table}[H]
\centering
\small
\begin{tabular}{p{0.27\textwidth}p{0.66\textwidth}}
\toprule
Policy & FCFS appointment assignment \\
Baseline demand & Class 1 = 50 arrivals/day; Class 2 = 50 arrivals/day \\
Capacity and horizon & 32 slots/day; 14-day booking horizon; 365 measured days \\
Varied in this section & #1 \\
Fixed at baseline & #2 \\
Grid and seeds & #3; #4 \\
\bottomrule
\end{tabular}
\end{table}
}

\newcommand{\mainfigure}[3]{%
\begin{figure}[H]
    \centering
    \includegraphics[width=\textwidth]{#1}
    \caption{#2}
    \label{#3}
\end{figure}
}

\title{FCFS Appointment Simulation: Readable Sensitivity Report}
\author{CUIMC Appointment Simulation}
\date{\today}

\begin{document}
\maketitle

\section{Executive Summary}

This report studies the first policy in the simulation: first-come, first-served (FCFS). The goal is to understand which parameters change total performance and which parameters benefit one class over the other.

\begin{table}[H]
\centering
\renewcommand{\arraystretch}{1.15}
\begin{tabular}{p{0.28\textwidth}p{0.64\textwidth}}
\toprule
Question & Short answer \\
\midrule
What matters most overall? & Total demand and no-show behavior. Higher demand lowers served rate and raises offered wait. Higher no-show risk wastes booked slots. \\
What creates class benefit? & Differences between the classes. The class with lower cancellation, lower balking, higher balking threshold, lower no-show risk, or higher no-show threshold tends to benefit. \\
Why not use utilization alone? & A busy clinic can still have poor access. In the baseline, utilization is high, but only about 27\% of arrivals complete service. \\
How should plots be read? & The line plots show simple one-at-a-time changes. The heatmaps show what happens when Class 1 and Class 2 parameters vary together. \\
\bottomrule
\end{tabular}
\end{table}

\section{Metrics and Baseline}

The analysis uses rates and averages, not total counts. This keeps scenarios comparable when arrival volume changes.

\begin{table}[H]
\centering
\renewcommand{\arraystretch}{1.12}
\begin{tabular}{p{0.24\textwidth}p{0.28\textwidth}p{0.39\textwidth}}
\toprule
Plain name & Code name & Meaning \\
\midrule
Served rate, Class $i$ & \texttt{percent\_serviced\_i} & Share of Class $i$ arrivals that complete service. \\
Overall served rate & \texttt{overall\_percent\_serviced} & Share of all arrivals that complete service. \\
Utilization & \texttt{average\_utilization} & Share of daily slots that become completed visits. This is aggregate only. \\
Offered wait & \texttt{mean\_offered\_booking\_delay} & Average delay offered to patients who received an offer, including patients who balked. \\
Class gap & \texttt{percent\_serviced\_1 - percent\_serviced\_2} & Positive means Class 1 has higher served rate. Negative means Class 2 has higher served rate. \\
\bottomrule
\end{tabular}
\end{table}

Baseline uses \texttt{configs/baseline.yaml}: $\lambda_1 = 50$, $\lambda_2 = 50$, 32 slots per day, 14-day horizon, balking threshold 9, balking high probability 0.50, no-show threshold 6, no-show high probability 0.30, and cancellation probability 0.10 for both classes.

\begin{table}[H]
\centering
\begin{tabular}{lrrrrrr}
\toprule
Scenario & Utilization & Overall served & Accepted wait & Offered wait & Class gap & Delay gap \\
\midrule
Baseline & 0.839 & 0.269 & 8.35 & 9.30 & 0.001 & 0.003 \\
Scenario 2 & 1.000 & 0.395 & 4.29 & 5.06 & -0.008 & 0.005 \\
\bottomrule
\end{tabular}
\caption{Scenario values are rates or averages.}
\end{table}

\section{How To Read The Figures}

Each behavior section has the same visual logic.

\begin{itemize}
    \item The slice plots are the easiest to read. One class changes while the other class stays at the baseline value.
    \item The interaction heatmap is for the two-variable question. It shows what happens when Class 1 and Class 2 values move together.
    \item Blue lines refer to Class 1, orange lines to Class 2, and black lines to the overall value.
    \item In class-gap heatmaps, blue means Class 1 has the higher served rate and orange means Class 2 has the higher served rate.
\end{itemize}

\section{Balking Step}

Balking step is the high-delay rejection probability after the balking threshold. A higher step means patients reject long-delay offers more often.

\simulationassumptions
{Class 1 and Class 2 balking step, from 0.00 to 1.00}
{Balking threshold = 9 days; no-show, cancellation, arrivals, capacity, horizon}
{21 values per class, step size 0.05}
{single fine-grid seed 2027}

\mainfigure{balking_step_slice_access.png}{Served-rate slices for balking step.}{fig:balk-step-access}
\mainfigure{balking_step_slice_utilization.png}{Utilization slices for balking step.}{fig:balk-step-util}
\mainfigure{balking_step_slice_wait.png}{Offered-wait slices for balking step.}{fig:balk-step-wait}
\mainfigure{balking_step_interaction_heatmap.png}{Two-variable balking-step heatmap. Left: overall served rate. Right: class gap.}{fig:balk-step-heatmap}

\textbf{Main reading.} Higher balking step hurts the class whose step increases because more of that class rejects long-delay offers. The other class can gain access to capacity that was released. Utilization moves much less than class served rates because total demand is still high.

\textbf{Takeaway.} Balking step is mostly redistribution. The class that balks less is more likely to be served.

\section{Balking Threshold}

Balking threshold is the delay where high balking begins. A lower threshold means patients start rejecting earlier. A higher threshold means patients tolerate longer waits.

\simulationassumptions
{Class 1 and Class 2 balking threshold, from 0 to 13 days}
{Balking step = 0.50; no-show, cancellation, arrivals, capacity, horizon}
{14 integer threshold values per class}
{single fine-grid seed 2027}

\mainfigure{balking_threshold_slice_access.png}{Served-rate slices for balking threshold.}{fig:balk-threshold-access}
\mainfigure{balking_threshold_slice_utilization.png}{Utilization slices for balking threshold.}{fig:balk-threshold-util}
\mainfigure{balking_threshold_slice_wait.png}{Offered-wait slices for balking threshold.}{fig:balk-threshold-wait}
\mainfigure{balking_threshold_interaction_heatmap.png}{Two-variable balking-threshold heatmap. Left: offered wait. Right: class gap.}{fig:balk-threshold-heatmap}

\textbf{Main reading.} The class with the higher threshold accepts longer waits and is more likely to complete service. Offered wait rises when thresholds are high because more patients stay in the appointment pool for long-delay offers.

\textbf{Takeaway.} Balking threshold is also mostly redistribution. The class that tolerates longer waits gets more service, but higher tolerance can keep the system more congested.

\section{No-Show Step}

No-show step is the high-delay probability that a booked patient does not attend. This happens after booking, so it mainly affects completed visits and slot use.

\simulationassumptions
{Class 1 and Class 2 no-show step, from 0.00 to 1.00}
{No-show threshold = 6 days; balking, cancellation, arrivals, capacity, horizon}
{21 values per class, step size 0.05}
{single fine-grid seed 2027}

\mainfigure{no_show_step_slice_access.png}{Served-rate slices for no-show step.}{fig:noshow-step-access}
\mainfigure{no_show_step_slice_utilization.png}{Utilization slices for no-show step.}{fig:noshow-step-util}
\mainfigure{no_show_step_slice_wait.png}{Offered-wait slices for no-show step.}{fig:noshow-step-wait}
\mainfigure{no_show_step_interaction_heatmap.png}{Two-variable no-show-step heatmap. Left: utilization. Right: class gap.}{fig:noshow-step-heatmap}

\textbf{Main reading.} Higher no-show step reduces completed visits. The class with higher no-show risk loses served rate after it has already accepted appointments. Offered wait changes little because the wait was assigned before the no-show event.

\textbf{Takeaway.} No-show step is both efficiency and redistribution. It wastes capacity and benefits the class with lower no-show risk.

\section{No-Show Threshold}

No-show threshold is the delay where high no-show risk begins. A higher threshold delays the point where no-show risk increases.

\simulationassumptions
{Class 1 and Class 2 no-show threshold, from 0 to 13 days}
{No-show step = 0.30; balking, cancellation, arrivals, capacity, horizon}
{14 integer threshold values per class}
{single fine-grid seed 2027}

\mainfigure{no_show_threshold_slice_access.png}{Served-rate slices for no-show threshold.}{fig:noshow-threshold-access}
\mainfigure{no_show_threshold_slice_utilization.png}{Utilization slices for no-show threshold.}{fig:noshow-threshold-util}
\mainfigure{no_show_threshold_slice_wait.png}{Offered-wait slices for no-show threshold.}{fig:noshow-threshold-wait}
\mainfigure{no_show_threshold_interaction_heatmap.png}{Two-variable no-show-threshold heatmap. Left: utilization. Right: class gap.}{fig:noshow-threshold-heatmap}

\textbf{Main reading.} A higher no-show threshold helps because patients become high-risk only at longer delays. The class with the higher threshold completes more booked appointments. Again, offered wait barely moves compared with served rate and utilization.

\textbf{Takeaway.} No-show threshold is mainly an efficiency parameter when both classes move together, and a redistribution parameter when thresholds differ by class.

\section{Cancellation Probability}

Cancellation is a direct probability that a booked appointment is canceled before the service day.

\simulationassumptions
{Class 1 and Class 2 cancellation probability, from 0.00 to 0.30}
{Balking, no-show, arrivals, capacity, horizon}
{21 values per class, step size 0.015}
{single fine-grid seed 2027}

\mainfigure{cancellation_probability_slice_access.png}{Served-rate slices for cancellation probability.}{fig:cancel-access}
\mainfigure{cancellation_probability_slice_utilization.png}{Utilization slices for cancellation probability.}{fig:cancel-util}
\mainfigure{cancellation_probability_slice_wait.png}{Offered-wait slices for cancellation probability.}{fig:cancel-wait}
\mainfigure{cancellation_probability_interaction_heatmap.png}{Two-variable cancellation heatmap. Left: overall served rate. Right: class gap.}{fig:cancel-heatmap}

\textbf{Main reading.} The class with higher cancellation probability loses completed visits. Canceled appointments may free capacity later, but the canceling patient is still not served.

\textbf{Takeaway.} Cancellation is one of the clearest class-gap drivers. The class with lower cancellation probability is much more likely to complete service.

\section{Arrival Rate and Class Mix}

Arrival parameters change how much demand enters the system. They do not directly change behavior, but they change congestion.

\simulationassumptions
{Total arrivals per day and Class 1 arrival share}
{Balking, no-show, cancellation, capacity, horizon}
{Total arrivals from 50 to 150 per day; Class 1 share from 0.10 to 0.90}
{single fine-grid seed 2027}

\mainfigure{arrival_rate_slice_access.png}{Served-rate slice as total arrivals change and class mix stays 50/50.}{fig:arrival-access}
\mainfigure{arrival_rate_slice_utilization.png}{Utilization slice as total arrivals change and class mix stays 50/50.}{fig:arrival-util}
\mainfigure{arrival_rate_slice_wait.png}{Offered-wait slice as total arrivals change and class mix stays 50/50.}{fig:arrival-wait}
\mainfigure{arrival_mix_interaction_heatmap.png}{Arrival-rate and class-mix heatmap. Left: overall served rate. Right: offered wait.}{fig:arrival-heatmap}

\textbf{Main reading.} Higher total demand lowers served rate and raises offered wait. Utilization can remain high even when access is poor. Changing the class mix alone does not create a strong class advantage because the two classes have the same behavior rules.

\textbf{Takeaway.} Arrival rate is mostly an efficiency and congestion parameter. Class mix matters more when the classes also differ in behavior.

\section{FCFS Stress Test}

The stress test keeps behavior fixed and varies total demand over a wider range. It shows why high utilization is not enough to judge performance.

\simulationassumptions
{Total arrivals per day, preserving the 50/50 class mix}
{Balking, no-show, cancellation, class mix, capacity, horizon}
{30 to 170 total arrivals per day}
{seeds 2027, 2028, 2029 averaged}

\mainfigure{fcfs_stress_access_wait.png}{FCFS stress test. Left: utilization versus served rate. Right: accepted and offered wait.}{fig:fcfs-stress}
\mainfigure{fcfs_arrival_outcome_decomposition.png}{Final outcome decomposition under FCFS. Each band is a share of arrivals.}{fig:fcfs-outcomes}

\begin{table}[H]
\centering
\resizebox{\textwidth}{!}{%
\begin{tabular}{rrrrrrrr}
\toprule
Arrivals/day & Utilization & Served & Offered wait & Balked & No offer & Lost after booking & Class gap \\
\midrule
30 & 0.943 & 0.998 & 0.22 & 0.000 & 0.000 & 0.002 & -0.000 \\
100 & 0.839 & 0.269 & 9.30 & 0.356 & 0.000 & 0.376 & 0.001 \\
170 & 0.841 & 0.157 & 10.59 & 0.242 & 0.346 & 0.255 & 0.001 \\
\bottomrule
\end{tabular}
}
\caption{Selected stress-test points. Outcome columns are shares of arrivals.}
\end{table}

\textbf{Main reading.} At low demand, almost every arrival is served. Around baseline demand, utilization is high but served rate is low. At overload, utilization is still high, but no-offers become a major loss channel.

\textbf{Takeaway.} FCFS can look busy while access is poor. Use served rate and outcome decomposition with utilization.

\section{Regression Screening}

The heatmaps change one or two parameters at a time. The regression screen asks which inputs explain outcomes across many random FCFS scenarios.

\begin{table}[H]
\centering
\small
\begin{tabular}{p{0.27\textwidth}p{0.66\textwidth}}
\toprule
Policy & FCFS appointment assignment \\
Randomized settings & 240 random parameter settings \\
Seeds & 2 seeds per setting, averaged \\
Inputs & Total demand, class mix, balking steps and thresholds, no-show steps and thresholds, cancellation probabilities \\
Outputs & Utilization, overall served rate, offered wait, Class 1 served-rate gap \\
\bottomrule
\end{tabular}
\end{table}

\begin{table}[H]
\centering
\begin{tabular}{llr}
\toprule
Target & Most important feature & Standardized coefficient \\
\midrule
Utilization & Average no-show threshold & 0.512 \\
Utilization & Average no-show step & -0.423 \\
Utilization & Average cancellation probability & 0.320 \\
Overall served rate & Total arrival rate & -0.774 \\
Overall served rate & Average no-show threshold & 0.251 \\
Offered wait & Total arrival rate & 0.576 \\
Offered wait & Average cancellation probability & -0.441 \\
Class gap & Cancellation probability gap & -0.452 \\
Class gap & Balking threshold gap & 0.429 \\
Class gap & Balking step gap & -0.349 \\
\bottomrule
\end{tabular}
\caption{Top regression drivers. Gap means Class 1 value minus Class 2 value.}
\end{table}

\textbf{Main reading.} Aggregate performance is mostly explained by total demand and average no-show behavior. Class advantage is mostly explained by differences between class behavior parameters.

\textbf{Takeaway.} Absolute demand and no-show levels explain overall performance. Class gaps explain who benefits.

\section{Final Takeaways}

\begin{enumerate}
    \item Served rate is the main access metric. Utilization alone can hide poor access.
    \item Balking mostly redistributes service toward the class that accepts longer waits.
    \item No-show behavior is a strong efficiency driver because booked slots can be wasted.
    \item Cancellation differences create large class gaps.
    \item Arrival rate controls congestion. Class mix alone is weak when classes behave the same.
    \item For class benefit, the difference between Class 1 and Class 2 parameters matters more than the absolute level.
\end{enumerate}

\appendix
\section{Appendix: Dense Reference Plots}

The main report uses readable slice plots and focused two-panel heatmaps. The figures below are retained for traceability when a fuller diagnostic view is useful.

\mainfigure{balking_threshold_jump_heatmaps.png}{Balking threshold and jump level when both classes share the same values.}{fig:appendix-balk-common}
\mainfigure{no_show_threshold_jump_heatmaps.png}{No-show threshold and jump level when both classes share the same values.}{fig:appendix-noshow-common}
\mainfigure{regression_standardized_coefficients.png}{Full standardized regression coefficient plot.}{fig:appendix-regression}
\mainfigure{fcfs_capacity_stress_curves.png}{Full FCFS stress diagnostic plot.}{fig:appendix-fcfs-full}

\end{document}
